\section{Introducción}
    El capturar un objeto del mundo real en imagen digital se podría considerar que es 
    todo un arte, debido a que no solamente se trata de contar con las herramientas 
    o sistemas adecuados para capturar (dígase cámaras u otros sistemas de imágenes), 
    sino también con las condiciones ambientales óptimas y las técnicas postproducción 
    para mejorar la calidad última de la imagen. Este proyecto trata de estos últimos 
    métodos, los cuales se van a emplear con el propósito de reducir el ruido visual 
    generado por un contraste o brillo inadecuados sobre toda la imagen o en canales 
    de colores específicos. \\

    Para la realización de este proyecto, se englobaron el uso de tres principales 
    técnicas: Histogramas de las imágenes, Corrección gamma y Transformación gain-bias. 
    El primero método permite mostrar la distribución de las intensidades (valores) 
    pertenecientes a los canales de colores de una imagen específicas, esto es auxiliar 
    para guiar qué transformación o tratamiento son los adecuados para mejorar visualmente 
    la imagen final, es decir, cambiar el contraste o brillo de forma conveniente.\\

    Mientras que las técnicas de corrección gamma y transformación gain-bias van a 
    permitir cambiar el contraste y brillo de una imagen de forma que sea la más 
    adecuada para el mejoramiento de la misma según un diagnóstico visual. La principal 
    distinción entre estas técnicas es que la corrección gamma opera de forma no 
    lineal mientras que la transformación gain-bias es lineal. También se considera 
    el uso de ecualización por histograma, aunque no se logró una mejora sustancial 
    en la calidad de las imágenes consideradas.